\documentclass[11pt,fleqn]{exam}
\usepackage[utf8]{inputenc}
\usepackage[T1]{fontenc}
\usepackage{fancyvrb}
\usepackage{amsmath}
\usepackage{amssymb}
\usepackage{hyperref}
\usepackage{algpseudocode}
\usepackage{comment}
\usepackage{enumitem}
\usepackage[margin=0.75in]{geometry}
\usepackage{tikz}
\usetikzlibrary{arrows,arrows.meta,positioning,intersections,shapes.gates.logic.US,calc}
\algnewcommand\algorithmicforeach{\textbf{for each}}
\algdef{S}[FOR]{ForEach}[1]{\algorithmicforeach\ #1\ \algorithmicdo}
\newcommand{\fillinMCmath}[1]{\begin{tikzpicture}\draw circle [radius=0.5em];\end{tikzpicture}\ #1}
\newcommand{\fillinMCmathsoln}[1]{\begin{tikzpicture}\draw[black, fill=blue] circle [radius=0.5em];\end{tikzpicture}\ #1}
\newcommand{\ptsamt}[1]{[#1~points]}
\newcommand{\ptamt}[1]{[#1~point]}

%%% Adding Colour to Questions and Answers
\usepackage{color}

\definecolor{solnblue}{rgb}{0,0,1}
\newenvironment{soln}{\color{solnblue}}{}

%Questions

% Answers
\definecolor{blu}{rgb}{0,0,0.5}
\def\blu#1{{\color{blu}#1}}
\definecolor{gre}{rgb}{0,.3,0}
\def\gre#1{{\color{gre}#1}}
\definecolor{red}{rgb}{0.5,0.0,0}
\def\red#1{{\color{red}#1}}
\def\norm#1{\|#1\|}
%%% End for Colours

\bracketedpoints

\makeatletter
\renewenvironment{solution}{\leavevmode\par\begin{soln}\noindent Solution:}{\end{soln}}
\makeatother
      \newif\ifsolutions\solutionsfalse

\author{}
\date{}
\title{CPSC 320 2025S: Assignment 1}

\begin{document}

\maketitle
\vspace{-0.5in} This assignment is due \textbf{Friday, July 11 at 7 PM}. Late submissions will not be accepted. Please follow these guidelines:
\begin{itemize}
	\item Prepare your  solution using \LaTeX \ and submit  a pdf file. Easiest will be to submit using
	      the .tex file provided. For questions where you  need to select a circle, you can simply
	      change \verb~\fillinMCmath~ to \verb~\fillinMCmathsoln~ .

	\item Enclose each paragraph of your solution with
	      \verb~\begin{soln}Your solution here...\end{soln}~.
	      \begin{soln}Your  solution will  then appear  in dark  blue\end{soln}, making  it a  lot
	      easier for TAs to find what you wrote.

	\item   Submit   the    assignment   via   \href{https://gradescope.ca/}{GradeScope}. You will be get access to Gradescope on the 9th of July via an invite link. Your group must make  a \textbf{single} submission via one
	      group member's account, marking all other group members in that submission \textbf{using
		      GradeScope's interface}.

	\item  After uploading  to  Gradescope, link  each  question  with the  page  of your  pdf
	      containing your solution.
\end{itemize}

Before we  begin, a few  notes on pseudocode throughout  CPSC 320: Your  pseudocode should
communicate your algorithm  clearly, concisely, correctly, and  without irrelevant detail.
Reasonable use  of plain  English is  fine in  such pseudocode.  You should  envision your
audience as a capable CPSC 320 student unfamiliar with the problem you are solving. You may \textbf{neither} include what we consider to be irrelevant coding details \textbf{nor} assume that  we understand the particular  language you
chose. (So, for example,  do not write \texttt{\#include <iostream>} at  the start of your
pseudocode,   and    avoid   language-specific   notation   like    C/C++/Java's   ternary
(question-mark-colon) operator.)

Remember also  to \textbf{justify/explain  your answers}. We  understand that  gauging how
much  justification to  provide can  be tricky.  Inevitably, judgment  is applied  by both
student and  grader as to how  much is enough, or  too much, and it's  frustrating for all
concerned  when judgments  don't align.  Justifications/explanations need  not be  long or
formal, but  should be  clear and  specific (referring  back to  lines of  pseudocode, for
example). Proofs should be a bit more formal.

On the  plus side, if  you choose an  incorrect answer when  selecting an option  but your
reasoning shows partial  understanding, you might get more marks  than if no justification
is provided. And  the effort you expend  in writing down the  justification will hopefully
help you gain deeper  understanding and may well help you converge  to the right selection
:).

\vspace{.1in}

Ok, time to get started...

\clearpage

\section*{Group Members}

Please list the CWLs of all group members here (even if you are submitting by yourself). We will deduct a mark if this is incorrect or missing.

- Jason Mac

\section{SMP Extreme True And/Or False}

Each of the following problems represents a SMP scenario (with $n$ employers and $n$ applicants) and a statement about that scenario. Each statement may be \textbf{always} true, \textbf{sometimes} true, or \textbf{never} true. Select the best of these three choices and then:
\begin{itemize}
	\item If the statement is \textbf{always} true, (a) give, and very briefly explain, an example instance in which it is true and (b) prove that it is always true.
	\item If the statement is \textbf{never} true, (a) give, and very briefly explain, an example instance in which it is false and (b) prove that it is always false.
	\item If the statement is \textbf{sometimes} true, (a) give, and very briefly explain, an example in which it is true and (b) give and very briefly explain an example instance in which it is false.
\end{itemize}
Here are the problems:

\begin{questions}
	\question[3] \textbf{Scenario:} an SMP instance with $n \ge 2$. \textbf{Statement:} there is exactly one stable matching.
	\ifsolutions\input{q1a-sol.tex}\fi

	\question[4] \textbf{Scenario:} an SMP instance with $n \ge 2$. \textbf{Statement:} there exists a stable matching in which nobody gets their first choice.
	\ifsolutions\input{q1b-sol.tex}\fi

	\question[4] \textbf{Scenario:} an SMP instance with $n \ge 2$ where all employers have identical preference lists. \textbf{Statement:} there is exactly one stable matching.
	\ifsolutions\input{q1c-sol.tex}\fi

\end{questions}

\ifsolutions\else\newpage\fi

\begin{soln}
	The statement is false. Take \(n = 2\) and we construct the following SMP instance of \(2\) employees and applicants denoted by
	\(E = \{e_1, e_2\}\) and \(A = \{a_1, a_2\}\) respectively. We define this instance through the following preference lists:

	\begin{itemize}
		\item \(a_1 : e_1, e_2\)
		\item \(a_2 : e_2, e_1\)
		\item \(e_1 : a_2, a_1\)
		\item \(e_2 : a_1, a_2\)
	\end{itemize}

	We have that a stable matching can be found through the pairings, \((e_1, a_1), (e_2, a_2)\), but also with \((e_1, a_2), (e_2, a_1)\).

	Now, for each \(n > 2\), we construct a new instance through the following:

	\begin{itemize}
		\item \(a_1 : e_1, e_2, \dots\)
		\item \(a_2 : e_2, e_1, \dots\)
		\item \(a_i : e_i, \dots\)
		\item \(e_1 : a_2, a_1, \dots\)
		\item \(e_2 : a_1, a_2, \dots\)
		\item \(e_i : a_i, \dots\)
	\end{itemize}

	for \(i \in \{3, 4, \dots, n\}\), and where the permutation of the preference list does not matter except for the specified entry.

	Then we have the two following stable matchings for such an instance:

	\(S_1 = \{(e_i, a_i) : i = 3, 4, \dots, n\} \cup \{(e_1, a_1), (e_2, a_2)\}\), and

	\(S_2 = \{(e_i, a_i) : i = 3, 4, \dots, n\} \cup \{(e_1, a_2), (e_2, a_1)\}\).

	We never have an instability since by construction, for \(3 \leq i \leq n\), \(e_i\) always prefers \(a_i\) over anyone else and vice versa.

	In other words, not every instance for \(n \geq 2\) will have exactly one stable matching, but there always exists at least one by \(G-S\) algorithm.



\end{soln}


\section{All-Stars}

Suppose you're organizing a tennis tournament between $n$ players which we simply label as $1,2, \ldots ,n$. Each player competes against every other player, and there are no repeat matches and no tied scores. We model the results of the competition as a directed graph with $n$ vertices and exactly one directed edge between each pair of vertices: $G=(V,E)$, with $V =\{1,2, \ldots ,n\}$ where either $(i,j)) \in E$ (if $i$ defeated $j$), or $(j,i)) \in E$ (if $j$ defeated $i$). Such a graph is called a \emph{tournament}.
		
You want to determine a way to rank these players, but a challenge here is that cycles \textbf{may} exist in your tournament -- e.g., $i$ may defeat $j$ in a match, and $j$ could defeat $k$, but then we could still have $k$ win a match against $i$; this  would lead to a cycle of length three. 

A ranking for the tournament is an ordering $(r_1,r_2, \ldots ,r_n)$ where node $r_1$ is established as the overall rank 1 player. Due to the potential existence of cycles it is challenging to find a fair ranking.


 \begin{questions}
    \question[4]  We decide that the top ranked player should 
 always be some node whose out-degree in the directed graph is maximum. We motivate this claim by proving that if player $i$ has maximum degree, then for every other player $j$ either $i$ beat $j$, or $i$ beat some player $k$ which beat $j$. Prove this claim. 
 
 {\em Hint: A proof by induction on $n$ works well.}

    \ifsolutions\begin{soln}
	Let \(G = (V, E)\) be a graph such that it is a tournament. Let \(i \in V\) be the node such that it has maximum out-degree.

	We aim to prove that \(\forall j \neq i \in V\) either \(i\) beat \(j\) or \(i\) beat \(k \in V\) who beat \(j\) for any size \(|V| = n \in \mathbb{N}\).

	Base case: \(|V| = 2\). Then if \(i\) has max out-degree then \(E = \{(i, j)\}\). Thus, the claim holds.

	Assume the claim holds for \(|V| = n\), with \(n \in \mathbb{N}\).

	Consider a tournament with \(|V| = n + 1\) nodes. Denote the player who has maximum out-degree by \(i\).

	Then remove \(j \neq i\) from this graph, so that we have \(n\) nodes, denote \(V' = V\setminus \{j\}\), with \(i\) beat \(j\).

	Either \(i\) has the maximum out-degree in this graph or not. We consider each case.

	\begin{itemize}
		\item \text{Case \(i\) has maximum out-degree in \(V'\)}. Then, by assumption for each node other node \(j' \in V'\), either \(i\) beat \(j'\) or \(i\) beat \(k'\) who beat \(j'\).
		      Since \(|V'| = n\).

		      By adding back in \(j\), then the statement holds for \(|V| = n + 1\), since \(i\) beat \(j\) by assumption.

		\item \text{Case \(i\) does not have maximum out-degree in \(V'\)}.

		      So, by removing \(j\), every other node decreased their out-degree by at most one.

		      By assumption \(i\) beat \(j\), so its out-degree must have decreased by one in \(V'\).

		      Thus, there was some node \(k \in V\) such that \(k\) had the same out-degree as \(i\), but \(k\) did beat \(j\).

		      Thus, \(j\) beat \(k\). By assumption for each other \(j' \in V'\), either \(k\) beat \(j'\) or \(k\) beat \(k'\) who beat \(j'\).

		      We now add back \(j\) to the graph. It remains that we need a connection to \(j\), as \(k\) did not beat \(j\).

		      Indeed if \(k\) beats \(i\) who beat \(j\), then assumption holds for \(|V| = n + 1\), using \(k\) in the hypothesis.

		      However, if \(k\) did not beat \(i\) then we consider removing \(i\) from this graph.

		      Since \(k\) lost to \(i\), then the out-degree of \(k\) cannot decrease. So \(k\) must have the max out-degree.

		      This graph being size \(n\), means there is some \(v \in V \setminus \{i\}\) so that \(k\) beats \(v\) and \(v\) beats \(j\).

		      Returning \(i\) to the graph, the statement holds for \(|V| = n + 1\), again using \(k\) in the hypothesis.
	\end{itemize}

	By induction, the statement holds for any tournament with \(|V| = n \in \mathbb{N}\).

\end{soln}
\fi
 
     \question[2] We decide that a reasonable definition  is suggested by the previous result. A sequence of nodes $(r_1,r_2, \ldots ,r_n)$ is a \emph{valid ranking} if for each $i=1,2, \ldots ,n$ the node $r_i$ has maximum out-degree in the sub-tournament induced by the nodes $r_i,r_{i+1}, \ldots ,r_n$. Give a polynomial-time algorithm which decides if a given permutation of the nodes $(v_1,v_2, \ldots ,v_n)$ is a valid ranking.

     \ifsolutions\begin{soln}
	Let \(P = (v_1, v_2, \dots, v_n)\) be a permutation on \(V\).

	We will assume that we have an adjency matrix to represent the edges in the graph \(M\).

	\begin{algorithmic}[1]
		\Procedure {Valid-Ranking}{P, M}
		\For{each $i = 1, 2, \dots, n - 1$}
		\State store the sum from $j = i$ to $j = n$ of $M(j)$ as $d_i$
		\For{each $j = i + 1, \dots, n - 1$}
		\State store the sum from $k = j$ to $k = n$ of $M(k)$ as $d_k$
		\If{$d_i < d_k$}
		\State end the procedure, report it is not valid
		\EndIf
		\EndFor
		\EndFor
		\State report the permutation as valid
		\EndProcedure
	\end{algorithmic}
	This algorithm is \(O(n^3)\), where \(n = |P|\). In the worse case, we have a valid permutation, and the algorithm does work as follows.

	We have to iterate through the permutation list \(n - 1\) times. Access to the elements will be \(O(1)\) as we assume we are given an array.

	Then for each iteration, we are iterating through the entire adjency matrix, except each iteration we are considering one less node.

	Access to the adjency matrix is \(O(1)\) thus the iteration of the adjency matrix is \(O(i^2)\).

	Then, for each \(i\), the work we are doing is \(i^2\). Thus, summing \(i = 1, 2, \dots, n - 1\) gives us \(O(n^3)\).


\end{soln}
\fi 

      \question[2] Describe a tournament of size $n$ which has precisely one valid ranking.

     \ifsolutions\begin{soln}
	Suppose we are given a tournament graph \(T = (V, E)\) with a vertex ordering \(v_1, v_2, \dots, v_n\) such that for all \(i < j\), the directed edge \((v_i, v_j)\) is in \(E\). That is, every vertex points to all vertices that come after it in the order.

	Suppose we have a tournament graph \(T = (V, E)\) such that there is an ordering, \(v_1, v_2, \dots, v_n\) so that there is an edge \((v_i, v_j)\) iff \(i < j\).

	This would create one valid ranking, and this ordering is exactly that ranking.

	Before we prove that statement, we will show it is a tournament in the first place, and that is is valid ranking.

	First, we require that every player has played someone else and only once or in other words, either \((v_i, v_j) \in E\) or \((v_j, v_i) \in E\) but not both.

	Let \(v_i, v_j \in V\) such that \(i \neq j\). Then either \(i < j\) or \(j < i\). Then we can only have one of the edge pairs in \(E\) by assumption.

	Next, we require that \(v_i\) has maximum out-degree in its subtournaments.

	Observe that for any \(v_i\) in a subtournament it's out degree is \(d_i = n - i + 1\), independent of the sub tournament.

	Then for any \(i < j\), we see that \(n - j + 1 < n - i + 1\), but this precisely means that \(d_j < d_i\), thus it is valid ranking.

	Proof: For contradiction, suppose that there were two valid rankings.

	Then it must be some permutation on \(P:= v_1, v_2, \dots, v_n\).

	Then there exists some \(i < j\) in \(P\) such that \(j < i\) in \(P'\). But this means that \(d_j < d_i\) in the subtournament for \(j\) in \(P'\).

	This contradicts that \(P'\) was a valid ranking.
\end{soln}
\fi 
      
     \question[4] One brute force method for  
      producing all valid rankings is to consider all permutations of the nodes and then use your algorithm for Part b to check each one if it is a valid ranking. This will run in time $O(n! p(n))$ where $p(n)$ is the runtime of your algorithm in Part b.

      Instead, create a better algorithm which runs in time $O(|Val| poly(n))$ where $Val$ is the set of valid rankings, and $poly(n)$ is a polynomial in $n$. (You do not need to analyze the runtime of your algorithm.)
 
     \ifsolutions\begin{soln}
	Find the set \(S_0\) such that each node in \(S_0\) has equal out-degree and is the greatest amongst all \(V\).

	Find the set \(S_1\) such that each node in \(S_1\) has equal out-degree and is the greatest amongst all \(V \setminus S_0\).

	Find set \(S_k\) so that each node in \(S_k\) has equal out-degree and is the greatest amongst all \(V \setminus \bigcup_{i=0}^{k-1}S_i\).

	Continue until you have a collection, such that \(\bigcup_{i = 0}^{k} S_i = V\).

	Then create a set of permutation, \(P_0, P_1, \dots, P_k\) such that \(P_i\) contains all possible permuations of ordered tuplets on vertices in \(S_i\).

	Next, create another set containing all ordered tuples of \(r = (p_1, p_2, \dots, p_k)\) where each \(p_i \in P_i\).

	Return \(R = \{\text{all possible } r\}\) as the set all possible valid rankings.

\end{soln}
\fi
 \end{questions}

\ifsolutions\else\newpage\fi



\begin{soln}
	This statement is never true.

	Let \(n \geq 2\), suppose we are given an instance an SMP. Denote set of employees by \(E\) and applicants by \(A\).
	We assume for the sake of contradiction, there exists a matching \(M\) such that nobody gets their first choice.
	Then in the G-S algorithm, this implies that for each \(e \in E\) they are either rejected on their first offer or
	their applicant to which they made their first offer to is poached by someone else.

	We also observe that the top choices of each employer is not unique, otherwise each employer would have never have been rejected in their first choice and our assumption is false.
	Then we have that there exists \(e_1, e_2 \in E\) such that their top choice of applicant is some \(a \in A\).
	Observe that the top employer for \(a\) cannot be \(e_1\) or \(e_2\) since this would mean that either \(e_1\) or \(e_2\) is never rejected by \(a\) or have \(a\) stolen from them.

	Thus, the top employer choice for \(a\) is some \(e \in E\).



\end{soln}

\section{Very Special Problems}

In this question, we consider special cases of NP-complete problems. Formally, we say that Problem B is a \textbf{special case} of Problem A if every instance of Problem B can be viewed as an instance of Problem A. We've seen examples of this in class: for example, 3-SAT is a special case of SAT (where every clause has length 3). Minimum spanning tree is a special case of Steiner Tree, in which the set of vertices we need to connect includes all the vertices in $V$.

\begin{questions}

	\question[4] Consider the \textbf{Bounded-Leaf Spanning Tree Problem (BLST)}: given a graph $G = (V,E)$ and an integer $k$, does $G$ have a spanning tree with no more than $k$ leaves?

	Give an NP-complete problem that is a special case of BLST, and justify why this problem is a special case.

	\ifsolutions\begin{soln}
	Let \(V\) be a vertex set. Add \(v \in V\). Then also add \(v_1, v_2, \dots, v_n\) to \(V\).

	Then for \(v_1, v_2, \dots, v_d\) add \((v, v_i) \in E\) for \(i = 1, 2, \dots, d\).

	Then colour the node \(v\) blue, and colour all other nodes in \(V\) red.

	By construction, the degree of all other nodes that are not \(v\) is \(1\) since it's only connected to \(v\).

	Thus, the degree of \(v\) is the maxmimum degree of the graph, which is \(d\).

	Since no nodes are adjacent except for ones to \(v\). Then the colour red is never shared by adjacent vertices.

	Thus, this graph with maximum degree \(d\) can be coloured with two colours.

\end{soln}
\fi

	\begin{soln}
		Let \(G = (V, E)\) be a graph. Suppose that this graph had a hamiltonian path.

		This means we can form a spanning tree out of this graph using exactly \(n - 1\) edges.

		Let \(P = (v_1, v_2, \dots, v_n)\) be the hamiltonian path in \(G\), and consider it the spanning tree structure.

		By definition, a leaf in a tree is a node with degree \(1\). We see that \(v_1, v_n\) both have degree \(1\).

		Hence, this tree has exactly two leaves.

		Thus, if we take \(G\) and ask the question to HAMPATH, does this graph contain a hamiltonian path?

		If the reply is yes, then \(G\) has a spanning tree with \(2\) leaves, or we can also say no more than \(2\) leaves.

		In other words, HAMPATH is a special case of BLST where we specify \(k = 2\).

	\end{soln}

	\question[3] You showed in the previous question that there is an NP-complete problem that is a special case of BLST, and it is not difficult to show that BLST is in NP (though we are not asking you to do this). Does this imply that BLST is NP-complete? Justify your answer.

	\ifsolutions\begin{soln}
	Proof: We provide a proof by induction on the number of nodes, \(n \in \mathbb{N}\), for graph \(G = (V, E)\).

	Base case \(|V| = 1\). The highest degree can be \(d = 0\), since the edge set would be empty.

	Thus, the algorithm uses \(1 = d + 1\) colours to colour the single node and the base case holds.

	Assume for any graph with \(n\) nodes that the algorithm will colour any ordering of nodes using at most \(d + 1\) colours.

	Consider a graph \(G = (V, E)\) with \(|V| = n + 1\) nodes with the highest degree of any node is \(d\).

	Let \(v_1, v_2, \dots, v_n, v_{n+1}\) be any ordering of the nodes. Remove \(v_{n+1}\) from this order, and the graph.

	Denote the deleted graph without \(v_{n+1}\) by \(G'\). Notice the degree of any node in \(G'\) can only decrease.

	Thus, the maximum degree of \(G'\) remains to be \(d\) with the removal of \(v_{n+1}\).

	By assumption, we can colour the ordering \(v_1, v_2, \dots, v_n\) using at most \(d+1\) colours.

	Since \(v_{n+1}\) has degree at most \(d\) then it has at most \(d\) neighbours. Then we add back \(v_{n+1}\).

	The number of colours used to colour its neigbours can be at most \(d\), if each are coloured distinctly.

	Thus, there remains at least \(1\) unique colour to colour \(v_{n+1}\) so that it is a valid colouring.

	In other words, the algorithm uses at most \(d + 1\) colours to colour the ordering \(v_1, v_2, \dots, v_{n+1}\).

	Induction makes the claim holds true for any graph with \(n\) nodes and maximum degree \(d\).


\end{soln}
\fi

	\begin{soln}
		This does not necessarily imply that BLST is NP-complete. Let \(Y\) be some NP-complete problem.

		By definition, to prove BLST is NP-complete, it remains to show that \(Y \leq_p BLST\).

		While we showed in the previous part that an NP-complete problem (HAMPATH) is a special case of BLST (for \( k = 2 \)), this is not sufficient to conclude NP-completeness of BLST in general.

		In particular, it remains to show that \(BLST(k = 2) \leq_p BLST\) or equivalently \(HAMPATH \leq_p BLST\).

		That is, we must show a polynomial time reduction from HAMPATH (or any another problem \(Y\)) to BLST, not just observe that one is a special case of the other.

		Hence, unless such a reduction is explicitly given with proof of correctness, we cannot conclude that BLST is NP-complete.
	\end{soln}


	\question[4] Give an example of a polytime-solvable problem you have seen in this class that is a special case of an NP-complete problem. Justify your answer.

	\ifsolutions\begin{soln}
	We consider the graph \(G = (V, E)\) with \(V = \{1, 2, 3, 4\}\) and \(E = \{(1, 4), (4, 3), (3, 2)\}\).

	\begin{center}
		\begin{tikzpicture}[node distance=2cm, auto]
			% Nodes
			\node[circle, draw, above  right = of 3] (1) {1};
			\node[circle, draw, right=of 1] (2) {2};
			\node[circle, draw, right=of 2] (3) {3};
			\node[circle, draw, above right=of 2] (4) {4};

			% Edges
			\draw (1) -- (4) -- (3) -- (2);
		\end{tikzpicture}
	\end{center}

	This graph can be coloured with three colours. Namely we can assign \(A = \{1, 3\}\) \(B = \{2, 4\}\).

	We see that no edges are shared between any vertices in \(A, B\) so this is a valid colouring.

	Now consider the ordering \(1, 2, 3, 4\). We first colour \(1\) blue. And then we consider \(2\).

	There is no edge between \(1, 2\) so we colour \(2\) blue and consider \(3\).

	So, its adjacent vertices have been coloured with blue, then we introduce a new colour for it red.

	Now we finally consider \(4\), its neighbours have been coloured with both red and green, so we must introduce a new colour for it purple.

	We have thus coloured this graph using three colours through the greedy algorithm when we could have used two.


\end{soln}
\fi

	\begin{soln}
		Recall the problem of finding a topological ordering on a directed graph \(G = (V, E)\).

		The dynamic programming solution can produce a solution in polynomial time.



		Now recall the
	\end{soln}

\end{questions}

\ifsolutions\else\newpage\fi

\section{With a Little Help from my Friends}

[10 points] On the popular social network Q, users ``follow'' other users to see the content they release. We can model this with a directed graph with a user $u_i$ represented by a vertex, and with a directed edge $(u_i, u_j)$ to indicate that $u_i$ follows $u_j$.

We say that two users $u_i$ and $u_j$ are \textbf{friends} if $u_i$ follows $u_j$ and $u_j$ follows $u_i$. We say that a group of users $W$ form a \textbf{friend group} if every user in $W$ is friends with every other user in $W$.

The \textbf{Friend Group Problem} (FGP) asks: given a directed graph $G = (U,E)$ representing the ``follows'' structure of our social network, does the network contain a friend group of size $k$?

Show how  to reduce an arbitrary  instance of FGP
into  an instance  of  \emph {Boolean  Satisfiability}  (SAT). To save time, we are not requiring you to prove the correctness of this reduction (though it's good practice if you want to give it a try), but you must clearly explain the purpose of all the components of your reduction. Hint:  my reduction uses $nk$ variables, where $n$ is the number of users in $U$.

\ifsolutions\input{q4-sol.tex}\fi

\ifsolutions\else\newpage\fi
\end{document}
